\begin{resumo}
As doenças cardiovasculares representam a principal causa de mortalidade no Brasil e no mundo, motivando o desenvolvimento de ferramentas computacionais de apoio ao diagnóstico clínico. Este trabalho compara dois paradigmas de redes neurais artificiais, o ADALINE Logístico e o Perceptron Multicamadas (MLP), na tarefa de predição de doenças cardíacas, avaliando cada arquitetura em versões de implementação própria (NumPy) e com biblioteca (scikit-learn), sob três configurações de treinamento sistematicamente variadas. Os experimentos foram conduzidos sobre dois conjuntos de dados públicos: o UCI Heart Disease Dataset (918 pacientes, triagem diagnóstica) e o Heart Failure Clinical Records Dataset (299 pacientes, prognóstico de mortalidade). Ao todo, 24 experimentos independentes foram avaliados não apenas pela acurácia global, mas principalmente pelo número de Falsos Negativos, critério de maior relevância clínica em diagnóstico médico. No Banco 1, o MLP Sklearn atingiu a maior acurácia absoluta (88,04\%), enquanto o MLP NumPy na configuração C3 (3.000 épocas, $\eta = 0{,}0001$) apresentou o menor número de Falsos Negativos (15), indicando superioridade clínica. No Banco 2, o MLP NumPy na configuração C1 (500 épocas, $\eta = 0{,}01$) ofereceu o melhor equilíbrio entre acurácia (78,89\%) e critério clínico (9 Falsos Negativos). Os resultados evidenciam que o ADALINE, modelo linearmente mais simples, mostrou-se competitivo com o MLP em diversas configurações, fenômeno explicado principalmente pelo volume limitado de dados disponíveis e, possivelmente, pela natureza predominantemente linear dos problemas em estudo, hipótese consistente com os resultados, mas empiricamente indistinguível da hipótese de insuficiência de dados. Um experimento complementar com os dados originais contínuos confirmou que a binarização das variáveis, adotada com o intuito de alinhar a representação às categorizações clínicas oficiais, não é o fator determinante da competitividade do ADALINE, e que os dados contínuos melhoram consistentemente os modelos ADALINE em ambos os bancos, com ganhos de até 13,3 pontos percentuais no Banco~2, enquanto os modelos MLP não se beneficiam de forma consistente da representação contínua.

\noindent Palavras-chave: ADALINE. Aprendizado de Máquina. Diagnóstico de Doenças Cardíacas. Falso Negativo. Perceptron Multicamadas. Predição de Doenças Cardíacas. Redes Neurais Artificiais.
\end{resumo}
